\chapter{Implementation and Software Design}
\label{ch:implementation}

\section{Chapter Overview}
\label{sec:imp_chapter_overview}
Turning theoretical computer vision models, planar homography geometry, and sequential deep learning into a working real-time virtual keyboard on standard plain paper requires a well-structured software architecture. Creating a camera-based human-computer interaction system that works with a single standard webcam, without specialized hardware (such as 3D depth sensors, Time-of-Flight cameras, laser projectors, or electronic gloves), requires clean software patterns, efficient multi-threading, accurate mathematical signal processing, and clear modular separation \cite{ReviewVirtualKeyboard2020, Maman2023TypeNet}.

Designing a monocular paper virtual keyboard presents several engineering challenges. Unlike physical mechanical keyboards or touchscreen devices, plain paper provides no physical switch movement, key travel, or electrical signals when a user presses a key. Therefore, touch detection must rely entirely on vision, tracking finger movements and speed changes across consecutive webcam frames \cite{Zhang2020MediaPipe, Ji2018Local}. In addition, a single regular RGB camera cannot directly measure 3D depth, creating ambiguity when a finger hovers close above the paper \cite{Thomas2013Camera}. Solving these issues requires combining vision algorithms, scale-invariant feature extraction, deep learning models, and desktop event dispatchers running on a standard CPU without requiring a GPU.

To make development and debugging manageable, the system components were built and tested across isolated partitions before assembling the complete system. The software suite consists of two main applications:
\begin{enumerate}
    \item \textbf{Application 1 (Layout Designer Suite \texttt{designer/designer\_app.py}):} An interactive PySide6 desktop GUI tool where users design custom button layouts, position keys, assign commands, and export synchronized printable PDF sheets (with embedded AprilTag markers for standard paper printing) and layout XML files.
    \item \textbf{Application 2 (Runtime Virtual Keyboard Engine and Command Mapper):} A runtime software application containing a Setup GUI (Component A) for mapping button IDs to keystrokes or shell commands across any standard webcam feed, and a Background Live Watch Engine (Component B) that tracks AprilTag markers on plain paper to compute the $3 \times 3$ Homography matrix, tracks MediaPipe hand landmarks, scale-normalizes joint coordinates by hand length ($L_{\text{hand}}$), evaluates 12~FPS 5-frame sliding windows using a PyTorch LSTM model (\texttt{best\_finger\_touch\_lstm.pth}) on a standard CPU, and triggers the mapped system actions.
\end{enumerate}

This chapter describes the implementation details of the software system. Section~\ref{sec:imp_framework_steps} explains the algorithm designs, workflow block diagrams, pseudocode algorithms, CPU timing budgets, and technology choices. Section~\ref{sec:imp_significant_attempts} details the implementation of the five research partitions and the two main application GUIs. Section~\ref{sec:imp_challenges} discusses practical engineering challenges faced during development and their software solutions. Finally, Section~\ref{sec:imp_summary} summarizes the chapter.

\section{Algorithmic Design and Software Engineering Steps}
\label{sec:imp_framework_steps}
Developing the virtual keyboard software followed an organized engineering process covering layout design, video capture, hand joint extraction, scale normalization, deep learning inference, AprilTag marker tracking, and system action execution.

\subsection{Framework Engineering Workflow Breakdown}
\label{subsec:imp_workflow_breakdown}
The complete system is split into two synchronized application tools working together across an end-to-end dataflow pipeline, as shown in the block diagram in Figure~\ref{fig:imp_block_diagram}.

\begin{figure}[htbp]
\centering
\resizebox{\textwidth}{!}{%
\begin{tikzpicture}[
    node distance = 0.8cm and 0.5cm,
    block/.style = {rectangle, draw=blue!80!black, fill=blue!10!white, thick, rounded corners=4pt, align=center, font=\sffamily\bfseries\tiny, minimum height=1.1cm, text width=2.4cm, inner sep=4pt},
    artblock/.style = {rectangle, draw=orange!80!black, fill=orange!10!white, thick, rounded corners=4pt, align=center, font=\sffamily\bfseries\tiny, minimum height=1.1cm, text width=2.4cm, inner sep=4pt},
    engineblock/.style = {rectangle, draw=green!80!black, fill=green!10!white, thick, rounded corners=4pt, align=center, font=\sffamily\bfseries\tiny, minimum height=1.1cm, text width=2.4cm, inner sep=4pt},
    arrow/.style = {draw=blue!80!black, ->, >=Stealth, thick}
]

% Column 1
\node (b1) [block] at (0, 0) {Application 1\\PySide6 Designer\\(\texttt{designer\_app.py})};

% Column 2
\node (b2) [artblock] at (3.5, 0) {Printable PDF Sheet\\(AprilTag Anchors)};
\node (b3) [artblock] at (3.5, -1.6) {Layout XML File\\(Bounding Boxes)};

% Column 3
\node (b4) [block] at (7.0, 0) {Application 2\\Setup \& Mapper GUI\\(Component A)};
\node (b5) [block] at (7.0, -1.6) {Camera Stream\\Selector \& Preview};
\node (b6) [block] at (7.0, -3.2) {Action Mapping\\Table (Keystrokes/Shell)};

% Column 4
\node (b7) [engineblock] at (10.5, 0) {Background Camera Thread\\(\texttt{camera\_thread.py})};
\node (b8) [engineblock] at (10.5, -1.6) {MediaPipe 21 Pose \&\\Unitless Scale Norm ($L_{\text{hand}}$)};
\node (b9) [engineblock] at (10.5, -3.2) {5-Frame Window Matrix\\($\mathbf{X}_W \in \mathbb{R}^{5 \times 84}$)};
\node (b10) [engineblock] at (10.5, -4.8) {PyTorch LSTM Classifier\\(Touch Model)};

% Column 5
\node (b11) [engineblock] at (14.0, -1.6) {AprilTag Tracker \& SVD\\Homography Engine ($H$)};
\node (b12) [engineblock] at (14.0, -4.8) {Planar Transformation\\($P_{\text{XML}} = H \cdot P_{\text{pixel}}$) \& Dispatch};

% Arrows
\draw [arrow] (b1) -- (b2);
\draw [arrow] (b1) -- (b3);
\draw [arrow] (b2) -- (b4);
\draw [arrow] (b3) -- (b4);
\draw [arrow] (b4) -- (b5);
\draw [arrow] (b5) -- (b6);
\draw [arrow] (b6) -- (b7);

\draw [arrow] (b7) -- (b8);
\draw [arrow] (b8) -- (b9);
\draw [arrow] (b9) -- (b10);
\draw [arrow] (b10) -- (b12);

\draw [arrow] (b7) -- (b11);
\draw [arrow] (b11) -- (b12);

\end{tikzpicture}%
}
\caption{End-to-End multi-module dataflow pipeline block diagram.}
\label{fig:imp_block_diagram}
\end{figure}

The dataflow pipeline runs through seven main processing stages:
\begin{enumerate}
    \item \textbf{Stage 1 (Layout Design and Export):} The user uses the PySide6 canvas in Application 1 (\texttt{designer\_app.py}) to create key buttons, snap them to an A4 grid, set button IDs, and define rectangular boundaries $[x_{\min}, y_{\min}, x_{\max}, y_{\max}]$. The exporter generates a vector PDF print sheet with corner AprilTag markers and an XML layout specification file.
    \item \textbf{Stage 2 (Action Mapping and Profile Saving):} The user opens Application 2 Component A (Setup GUI). The software loads the XML layout file, reads the key IDs, and shows an action mapping table. The user assigns each key ID to a single key (such as `'A'`), shortcut (such as `'Ctrl+C'`), or shell command. The physical layout is separated from digital actions: users can save and switch different action profiles (JSON files) for the exact same printed paper sheet without reprinting.
    \item \textbf{Stage 3 (Video Frame Ingestion):} Application 2 Component B starts running in the background. A dedicated worker thread (\texttt{camera\_thread.py}) captures webcam frames using OpenCV \texttt{VideoCapture} and puts raw frame arrays $\mathbf{I}_{\text{frame}} \in \mathbb{R}^{H \times W \times 3}$ into a thread-safe Queue buffer.
    \item \textbf{Stage 4 (AprilTag Tracking and Homography Calculation):} The tracking thread reads frames from the queue, converts them to grayscale, detects AprilTag quads (\texttt{tag36h11} family), extracts sub-pixel corner points, builds the Direct Linear Transformation (DLT) matrix $\mathbf{A}$, and calculates the Planar Homography matrix $H \in \mathbb{R}^{3 \times 3}$ using Singular Value Decomposition (SVD).
    \item \textbf{Stage 5 (Hand Joint Extraction and Scale Normalization):} Concurrently, MediaPipe Hand Landmarker extracts 21 2D hand joint coordinates $\mathbf{P}_i(t) = (x_i(t), y_i(t))$. Coordinates are shifted relative to wrist Landmark 0 ($\mathbf{P}_0$) and scaled by hand length $L_{\text{hand}}(t) = \|\mathbf{P}_9(t) - \mathbf{P}_0(t)\|_2$. Velocity vectors $\mathbf{V}_i(t)$ are calculated across consecutive frames ($\Delta t = \frac{1}{12}$~s). Normalized values are passed directly to the next stage to preserve sudden deceleration changes.
    \item \textbf{Stage 6 (Sliding Window and PyTorch LSTM Inference):} The feature pipeline collects 84-dimensional feature vectors $\mathbf{f}_t$ across 5 consecutive frames into temporal matrix $\mathbf{X}_W \in \mathbb{R}^{5 \times 84}$. The PyTorch model (\texttt{best\_finger\_touch\_lstm.pth}) evaluates $\mathbf{X}_W$ in under $2.1$~ms on a standard CPU, predicting touch probabilities $\mathbf{P}_{\text{touch}} \in [0, 1]^5$. If probability $p_k > 0.90$, finger $k$ is recognized as touching the paper.
    \item \textbf{Stage 7 (Planar Coordinate Mapping and Action Triggering):} When a touch is confirmed, active fingertip pixel coordinates $(u_{\text{active}}, v_{\text{active}})$ are mapped through $H$ into layout coordinates $(X_{\text{target}}, Y_{\text{target}}) = H \cdot [u, v, 1]^T / w$. The system checks which button bounding box contains this coordinate and triggers the corresponding keypress or shell command.
\end{enumerate}

\subsection{Mathematical Algorithm Formulations and Formal Pseudocode}
\label{subsec:imp_pseudocode_algorithms}
To clearly describe the core software operations, mathematical formulations and formal pseudocode algorithms are presented for the main processing engines.

\subsubsection{Algorithm 1: AprilTag Corner Localization and SVD Planar Homography Engine}
The homography engine detects four AprilTag marker corners and computes the $3 \times 3$ Planar Homography matrix $H$ that maps camera pixel coordinates $(u, v)$ to layout XML coordinates $(x, y)$.

Given $N \ge 4$ point correspondences $(u_i, v_i) \leftrightarrow (x_i, y_i)$, the Direct Linear Transformation (DLT) forms the linear equation $\mathbf{A}_{2N \times 9} \mathbf{h} = \mathbf{0}$, where:
\begin{equation}
\mathbf{A}_i = \begin{bmatrix}
-u_i & -v_i & -1 & 0 & 0 & 0 & u_i x_i & v_i x_i & x_i \\
0 & 0 & 0 & -u_i & -v_i & -1 & u_i y_i & v_i y_i & y_i
\end{bmatrix}
\label{eq:dlt_matrix_app}
\end{equation}

Solving $\mathbf{A} = \mathbf{U} \mathbf{\Sigma} \mathbf{V}^T$ using Singular Value Decomposition (SVD), the column vector $\mathbf{h}$ corresponding to the smallest singular value in $\mathbf{V}$ forms the homography matrix $H \in \mathbb{R}^{3 \times 3}$.

Algorithm~\ref{alg:homography_svd} presents the formal pseudocode for AprilTag corner localization and SVD homography calculation.

\begin{figure}[htbp]
\centering
\begin{minipage}{0.92\textwidth}
\small
\setlength{\extrarowheight}{2pt}
\begin{tabular}{|p{0.96\linewidth}|}
\hline
\textbf{Algorithm 1: AprilTag Corner Localization and SVD Planar Homography Calculation} \\
\hline
{\small\sffamily \textbf{Input:} Video frame image $\mathbf{I}_{\text{frame}} \in \mathbb{R}^{H \times W \times 3}$, Target XML anchor coordinates $\{(x_k^{\text{xml}}, y_k^{\text{xml}})\}_{k=1}^M$} \\
{\small\sffamily \textbf{Output:} Updated Planar Homography matrix $H \in \mathbb{R}^{3 \times 3}$, Sub-pixel tracking error $e_{\text{rms}}$} \\
{\small 1: Convert $\mathbf{I}_{\text{frame}}$ to grayscale single-channel image $\mathbf{I}_{\text{gray}}$.} \\
{\small 2: Detect AprilTag visual quads using AprilTag family \texttt{tag36h11}.} \\
{\small 3: \textbf{if} Count of detected valid tags $M_{\text{detected}} < 4$ \textbf{then}} \\
{\small 4: \hspace{1em} Retain previous valid Homography matrix $H_{\text{prev}}$ from memory buffer.} \\
{\small 5: \hspace{1em} \textbf{return} $H_{\text{prev}}$, Status: \texttt{TAG\_OCCLUSION\_WARNING}} \\
{\small 6: \textbf{end if}} \\
{\small 7: Extract sub-pixel corner pixel coordinates $\{(u_k, v_k)\}_{k=1}^{4 M_{\text{detected}}}$.} \\
{\small 8: Construct $2N \times 9$ DLT coefficient matrix $\mathbf{A}$ using Equation~\ref{eq:dlt_matrix_app}.} \\
{\small 9: Compute Singular Value Decomposition: $\mathbf{U}, \mathbf{\Sigma}, \mathbf{V}^T = \text{SVD}(\mathbf{A})$.} \\
{\small 10: Extract column vector $\mathbf{h} = \mathbf{V}[:, 8]$ corresponding to minimum singular value $\sigma_{\min}$.} \\
{\small 11: Reshape $\mathbf{h}$ into $3 \times 3$ matrix $H_{\text{raw}}$ and normalize $H = \frac{1}{h_{33}} H_{\text{raw}}$.} \\
{\small 12: Compute RMS reprojection error $e_{\text{rms}} = \sqrt{\frac{1}{N} \sum_{i=1}^N \| \mathbf{P}_i^{\text{xml}} - H \cdot \mathbf{P}_i^{\text{pixel}} \|^2}$.} \\
{\small 13: Update shared runtime memory buffer with $H$.} \\
{\small 14: \textbf{return} $H, e_{\text{rms}}$} \\
\hline
\end{tabular}
\end{minipage}
\caption{Formal algorithmic specification for AprilTag corner localization and SVD planar homography calculation.}
\label{alg:homography_svd}
\end{figure}

\subsubsection{Algorithm 2: Unitless Hand Landmark Scale Normalization Engine}
The pose engine extracts 21 anatomical 2D hand joint coordinates $\mathbf{P}_i(t) = (x_i(t), y_i(t))$ per frame and applies scale normalization.

Wrist Landmark 0 is used as origin $\mathbf{P}_0(t)$. Hand length $L_{\text{hand}}(t)$ is measured between wrist Landmark 0 and middle MCP Landmark 9:
\begin{equation}
L_{\text{hand}}(t) = \sqrt{(x_9(t) - x_0(t))^2 + (y_9(t) - y_0(t))^2}
\label{eq:hand_length_app}
\end{equation}

Unitless scale-normalized coordinates $\mathbf{P}_i^{\text{norm}}(t)$ are calculated as follows:
\begin{equation}
\mathbf{P}_i^{\text{norm}}(t) = \left( \frac{x_i(t) - x_0(t)}{L_{\text{hand}}(t)}, \frac{y_i(t) - y_0(t)}{L_{\text{hand}}(t)} \right)
\label{eq:norm_coords_app}
\end{equation}

Spatial velocity vectors $\mathbf{V}_i(t) = (v_{x,i}(t), v_{y,i}(t))$ between consecutive frames are derived at $\Delta t = \frac{1}{12}$~s:
\begin{equation}
v_{x,i}(t) = \frac{x_i^{\text{norm}}(t) - x_i^{\text{norm}}(t-1)}{\Delta t}, \quad
v_{y,i}(t) = \frac{y_i^{\text{norm}}(t) - y_i^{\text{norm}}(t-1)}{\Delta t}
\label{eq:velocity_app}
\end{equation}

These scale-normalized coordinates and velocities are passed directly to feature windowing to preserve deceleration signals when a finger hits the paper.

Algorithm~\ref{alg:scale_normalization} details the formal pseudocode for hand landmark scale normalization and velocity calculation.

\begin{figure}[htbp]
\centering
\begin{minipage}{0.92\textwidth}
\small
\setlength{\extrarowheight}{2pt}
\begin{tabular}{|p{0.96\linewidth}|}
\hline
\textbf{Algorithm 2: Unitless Hand Landmark Scale Normalization and Velocity Derivation} \\
\hline
{\small\sffamily \textbf{Input:} Raw MediaPipe 21 landmark 2D coordinates $\{\mathbf{P}_i(t) = (x_i, y_i)\}_{i=0}^{20}$, Previous frame keypoints $\{\mathbf{P}_i(t-1)\}_{i=0}^{20}$, Frame time step $\Delta t = \frac{1}{12}\text{ s}$} \\
{\small\sffamily \textbf{Output:} 84-Dimensional per-frame feature vector $\mathbf{f}_t \in \mathbb{R}^{84}$} \\
{\small 1: Extract wrist origin coordinates $\mathbf{P}_0(t) = (x_0(t), y_0(t))$.} \\
{\small 2: Extract middle MCP landmark coordinates $\mathbf{P}_9(t) = (x_9(t), y_9(t))$.} \\
{\small 3: Compute unitless hand length scale factor: $L_{\text{hand}}(t) = \sqrt{(x_9(t) - x_0(t))^2 + (y_9(t) - y_0(t))^2}$.} \\
{\small 4: \textbf{if} $L_{\text{hand}}(t) < \epsilon$ ($< 10$ pixels) \textbf{then}} \\
{\small 5: \hspace{1em} \textbf{return} Status: \texttt{INVALID\_HAND\_SCALE\_ERROR}} \\
{\small 6: \textbf{end if}} \\
{\small 7: Initialize empty feature vector array $\mathbf{f}_t \in \mathbb{R}^{84}$.} \\
{\small 8: \textbf{for} each landmark $i \in \{0, 1, \dots, 20\}$ \textbf{do}} \\
{\small 9: \hspace{1em} Calculate normalized position: $x_i^{\text{norm}}(t) = \frac{x_i(t) - x_0(t)}{L_{\text{hand}}(t)}$, \hspace{1em} $y_i^{\text{norm}}(t) = \frac{y_i(t) - y_0(t)}{L_{\text{hand}}(t)}$.} \\
{\small 10: \hspace{1em} Calculate numerical velocity: $v_{x,i}(t) = \frac{x_i^{\text{norm}}(t) - x_i^{\text{norm}}(t-1)}{\Delta t}$, \hspace{1em} $v_{y,i}(t) = \frac{y_i^{\text{norm}}(t) - y_i^{\text{norm}}(t-1)}{\Delta t}$.} \\
{\small 11: \hspace{1em} Store $(x_i^{\text{norm}}, y_i^{\text{norm}}, v_{x,i}, v_{y,i})$ into feature vector $\mathbf{f}_t$ at indices $[4i : 4i+4]$.} \\
{\small 12: \textbf{end for}} \\
{\small 13: Direct raw feature propagation (preserve contact deceleration transients).} \\
{\small 14: \textbf{return} $\mathbf{f}_t \in \mathbb{R}^{84}$} \\
\hline
\end{tabular}
\end{minipage}
\caption{Formal algorithmic specification for unitless hand landmark scale normalization.}
\label{alg:scale_normalization}
\end{figure}

\subsubsection{Algorithm 3: 5-Frame Sliding Window Assembly and PyTorch LSTM Touch Inference}
The temporal inference engine stacks 84-dimensional feature vectors $\mathbf{f}_t$ across 5 consecutive frames into input matrix $\mathbf{X}_W \in \mathbb{R}^{5 \times 84}$.

The PyTorch uni-directional LSTM evaluates input matrix $\mathbf{X}_W$:
\begin{equation}
\mathbf{h}_t, \mathbf{c}_t = \text{LSTM}(\mathbf{X}_W, \mathbf{h}_{t-1}, \mathbf{c}_{t-1})
\label{eq:lstm_forward_app}
\end{equation}

A fully connected dense layer with sigmoid activation predicts the binary touch probability vector $\mathbf{P}_{\text{touch}} = [p_{\text{thumb}}, p_{\text{index}}, p_{\text{middle}}, p_{\text{ring}}, p_{\text{pinky}}] \in [0, 1]^5$. If probability $p_k > 0.90$, digit $k$ is classified as touching the paper.

Algorithm~\ref{alg:lstm_touch_inference} details the formal pseudocode for 5-frame sliding window assembly and PyTorch LSTM touch inference.

\begin{figure}[htbp]
\centering
\begin{minipage}{0.92\textwidth}
\small
\setlength{\extrarowheight}{2pt}
\begin{tabular}{|p{0.96\linewidth}|}
\hline
\textbf{Algorithm 3: 5-Frame Sliding Window Assembly and PyTorch LSTM Touch Inference} \\
\hline
{\small\sffamily \textbf{Input:} Per-frame feature vector $\mathbf{f}_t \in \mathbb{R}^{84}$, Sliding window Queue buffer $\mathbf{Q}_{\text{win}}$, PyTorch model instance \texttt{best\allowbreak\_finger\allowbreak\_touch\allowbreak\_lstm.pth}, Detection threshold $\tau_{\text{touch}} = 0.90$} \\
{\small\sffamily \textbf{Output:} Active finger touch status array $\mathbf{S}_{\text{touch}} \in \{0, 1\}^5$, Active fingertip landmark index $k_{\text{active}}$} \\
{\small 1: Append current feature vector $\mathbf{f}_t$ to sliding queue $\mathbf{Q}_{\text{win}}$.} \\
{\small 2: \textbf{if} Length of $\mathbf{Q}_{\text{win}} < 5$ frames \textbf{then}} \\
{\small 3: \hspace{1em} \textbf{return} $\mathbf{S}_{\text{touch}} = [0,0,0,0,0]$, Status: \texttt{WINDOW\_BUFFERING}} \\
{\small 4: \textbf{end if}} \\
{\small 5: Construct sliding window input matrix $\mathbf{X}_W = [\mathbf{f}_{t-4}, \mathbf{f}_{t-3}, \mathbf{f}_{t-2}, \mathbf{f}_{t-1}, \mathbf{f}_t]^T \in \mathbb{R}^{5 \times 84}$.} \\
{\small 6: Advance sliding queue by stride of 3 frames (retaining 2-frame overlap).} \\
{\small 7: Convert $\mathbf{X}_W$ to PyTorch Tensor: $\mathbf{T}_{\text{input}} = \text{torch.tensor}(\mathbf{X}_W, \text{dtype}=\text{float32}).\text{unsqueeze}(0)$.} \\
{\small 8: Execute PyTorch model forward pass in evaluation mode (\texttt{with torch.no\_grad():}):} \\
{\small 9: \hspace{1em} $\mathbf{y}_{\text{logits}} = \text{Model}(\mathbf{T}_{\text{input}})$.} \\
{\small 10: \hspace{1em} $\mathbf{P}_{\text{touch}} = \text{sigmoid}(\mathbf{y}_{\text{logits}}) \in [0, 1]^5$.} \\
{\small 11: Initialize binary status vector $\mathbf{S}_{\text{touch}} = [0, 0, 0, 0, 0]$.} \\
{\small 12: \textbf{for} each finger digit $k \in \{0, 1, 2, 3, 4\}$ \textbf{do}} \\
{\small 13: \hspace{1em} \textbf{if} $P_{\text{touch}}[k] > \tau_{\text{touch}}$ \textbf{then}} \\
{\small 14: \hspace{2em} Set $\mathbf{S}_{\text{touch}}[k] = 1$ (Touch Confirmed).} \\
{\small 15: \hspace{2em} Set active fingertip landmark index: $k_{\text{active}} \in \{4, 8, 12, 16, 20\}$.} \\
{\small 16: \hspace{1em} \textbf{end if}} \\
{\small 17: \textbf{end for}} \\
{\small 18: \textbf{return} $\mathbf{S}_{\text{touch}}, k_{\text{active}}$} \\
\hline
\end{tabular}
\end{minipage}
\caption{Formal algorithmic specification for 5-frame sliding window assembly and PyTorch LSTM touch inference.}
\label{alg:lstm_touch_inference}
\end{figure}

\subsubsection{Algorithm 4: Planar Coordinate Mapping and Action Execution Engine}
When Algorithm~\ref{alg:lstm_touch_inference} detects a touch event, Algorithm~\ref{alg:coordinate_mapping} maps the active fingertip pixel coordinates $(u_{\text{active}}, v_{\text{active}})$ through Homography matrix $H$ into layout XML coordinates:
\begin{equation}
\begin{bmatrix} x_{\text{xml}} \\ y_{\text{xml}} \\ w \end{bmatrix} = H \cdot \begin{bmatrix} u_{\text{active}} \\ v_{\text{active}} \\ 1 \end{bmatrix}, \quad
X_{\text{target}} = \frac{x_{\text{xml}}}{w}, \quad Y_{\text{target}} = \frac{y_{\text{xml}}}{w}
\label{eq:homog_mapping_app}
\end{equation}

The system searches the XML key bounding boxes. If $(X_{\text{target}}, Y_{\text{target}}) \in [x_{\min}^j, x_{\max}^j] \times [y_{\min}^j, y_{\max}^j]$, key $j$ is pressed, executing the assigned keypress or shell command.

Algorithm~\ref{alg:coordinate_mapping} shows the formal pseudocode for planar coordinate mapping and action execution.

\begin{figure}[htbp]
\centering
\begin{minipage}{0.92\textwidth}
\small
\setlength{\extrarowheight}{2pt}
\begin{tabular}{|p{0.96\linewidth}|}
\hline
\textbf{Algorithm 4: Planar Coordinate Mapping and Action Execution Engine} \\
\hline
{\small\sffamily \textbf{Input:} Active fingertip pixel coordinate $\mathbf{P}_{\text{pixel}} = (u, v)$, Homography matrix $H \in \mathbb{R}^{3 \times 3}$, Target XML key bounding boxes $\{B_j = (x_{\min}^j, y_{\min}^j, x_{\max}^j, y_{\max}^j)\}_{j=1}^K$, Action mapping dictionary $\mathbf{M}_{\text{action}}$} \\
{\small\sffamily \textbf{Output:} Executed Action Status, Hit Key ID $j_{\text{hit}}$} \\
{\small 1: Compute homogeneous target vector: $[x_{\text{temp}}, y_{\text{temp}}, w]^T = H \cdot [u, v, 1]^T$.} \\
{\small 2: \textbf{if} $|w| < 10^{-6}$ \textbf{then}} \\
{\small 3: \hspace{1em} \textbf{return} Status: \texttt{HOMOGRAPHY\_SINGULARITY\_ERROR}} \\
{\small 4: \textbf{end if}} \\
{\small 5: Dehomogenize XML planar coordinates: $X_{\text{xml}} = \frac{x_{\text{temp}}}{w}$, \hspace{1em} $Y_{\text{xml}} = \frac{y_{\text{temp}}}{w}$.} \\
{\small 6: Initialize $j_{\text{hit}} = \text{None}$.} \\
{\small 7: \textbf{for} each key button $j \in \{1, 2, \dots, K\}$ in XML layout \textbf{do}} \\
{\small 8: \hspace{1em} \textbf{if} $x_{\min}^j \le X_{\text{xml}} \le x_{\max}^j$ \textbf{and} $y_{\min}^j \le Y_{\text{xml}} \le y_{\max}^j$ \textbf{then}} \\
{\small 9: \hspace{2em} Set $j_{\text{hit}} = j$.} \\
{\small 10: \hspace{2em} Extract action binding: $A_j = \mathbf{M}_{\text{action}}[j]$.} \\
{\small 11: \hspace{2em} \textbf{if} $A_j$ is \texttt{KeystrokeAction} \textbf{then}} \\
{\small 12: \hspace{2em} \hspace{1em} Inject system virtual key press event via OS API (e.g., `'Ctrl+C'`).} \\
{\small 13: \hspace{2em} \textbf{else if} $A_j$ is \texttt{SystemCommandAction} \textbf{then}} \\
{\small 14: \hspace{2em} \hspace{1em} Execute background shell command script: \texttt{subprocess.Popen(A\_j.command)}.} \\
{\small 15: \hspace{2em} \textbf{end if}} \\
{\small 16: \hspace{2em} Break loop.} \\
{\small 17: \hspace{1em} \textbf{end if}} \\
{\small 18: \textbf{end for}} \\
{\small 19: \textbf{return} Status: \texttt{ACTION\_DISPATCHED\_SUCCESS}, Hit Key ID $j_{\text{hit}}$} \\
\hline
\end{tabular}
\end{minipage}
\caption{Formal algorithmic specification for planar coordinate mapping and action execution.}
\label{alg:coordinate_mapping}
\end{figure}

\subsubsection{Algorithm 5: PySide6 Interactive Canvas Viewport Drag-and-Drop Math}
Application 1 (\texttt{designer\_app.py}) provides interactive mouse drag-and-drop key placement on a 2D canvas.

Algorithm~\ref{alg:canvas_drag_drop} details formal pseudocode for mouse viewport coordinate mapping and snap-to-grid calculations.

\begin{figure}[htbp]
\centering
\begin{minipage}{0.92\textwidth}
\small
\setlength{\extrarowheight}{2pt}
\begin{tabular}{|p{0.96\linewidth}|}
\hline
\textbf{Algorithm 5: PySide6 Interactive Canvas Viewport Drag-and-Drop Math} \\
\hline
{\small\sffamily \textbf{Input:} Viewport mouse click coordinate $\mathbf{p}_{\text{mouse}} = (u_{\text{m}}, v_{\text{m}})$, Viewport transformation matrix $\mathbf{T}_{\text{view}} \in \mathbb{R}^{3 \times 3}$, Grid snap spacing $\Delta g = 5.0\text{ mm}$, Button size $(W_{\text{key}}, H_{\text{key}})$} \\
{\small\sffamily \textbf{Output:} Snapped Key Bounding Box $B_{\text{new}} = (x_{\min}, y_{\min}, x_{\max}, y_{\max})$} \\
{\small 1: Map viewport pixel coordinate to scene millimeter space: $[x_{\text{scene}}, y_{\text{scene}}, 1]^T = \mathbf{T}_{\text{view}}^{-1} \cdot [u_{\text{m}}, v_{\text{m}}, 1]^T$.} \\
{\small 2: Compute grid-snapped origin coordinates:} \\
{\small 3: \hspace{1em} $x_{\text{snap}} = \text{round}(x_{\text{scene}} / \Delta g) \times \Delta g$.} \\
{\small 4: \hspace{1em} $y_{\text{snap}} = \text{round}(y_{\text{scene}} / \Delta g) \times \Delta g$.} \\
{\small 5: Define bounding box: $x_{\min} = x_{\text{snap}}$, $y_{\min} = y_{\text{snap}}$, $x_{\max} = x_{\text{snap}} + W_{\text{key}}$, $y_{\max} = y_{\text{snap}} + H_{\text{key}}$.} \\
{\small 6: Check bounding box collision against existing keys: $\text{Overlap} = \text{CheckIntersections}(B_{\text{new}}, \{B_{\text{existing}}\})$.} \\
{\small 7: \textbf{if} Overlap is True \textbf{then}} \\
{\small 8: \hspace{1em} Shift origin $x_{\text{snap}} = x_{\text{snap}} + \Delta g$ and recompute $B_{\text{new}}$.} \\
{\small 9: \textbf{end if}} \\
{\small 10: Instantiate \texttt{KeyButtonGraphic} graphics item on \texttt{QGraphicsScene}.} \\
{\small 11: \textbf{return} $B_{\text{new}}$} \\
\hline
\end{tabular}
\end{minipage}
\caption{Formal algorithmic specification for PySide6 canvas viewport drag-and-drop key placement.}
\label{alg:canvas_drag_drop}
\end{figure}

\subsubsection{Algorithm 6: Multi-Threaded Consumer-Producer Queue Synchronization}
Application 2 (\texttt{camera\_thread.py}) uses multi-threaded queues to separate camera capture from model inference.

Algorithm~\ref{alg:queue_synchronization} details formal pseudocode for consumer-producer thread synchronization.

\begin{figure}[htbp]
\centering
\begin{minipage}{0.92\textwidth}
\small
\setlength{\extrarowheight}{2pt}
\begin{tabular}{|p{0.96\linewidth}|}
\hline
\textbf{Algorithm 6: Multi-Threaded Consumer-Producer Queue Synchronization} \\
\hline
{\small\sffamily \textbf{Input:} Video Capture object \texttt{cv2.VideoCapture}, Shared Queue buffer $\mathbf{Q}_{\text{frame}}$ (Max Size = 2), Mutex Lock $\mathbf{L}_{\text{mutex}}$} \\
{\small\sffamily \textbf{Output:} Thread-safe Frame Stream Buffer} \\
{\small 1: \textbf{Worker Thread Function} \texttt{CameraCaptureLoop()}:} \\
{\small 2: \textbf{while} Runtime flag \texttt{is\_running} is True \textbf{do}} \\
{\small 3: \hspace{1em} Read raw video frame: $\text{ret}, \mathbf{I}_{\text{raw}} = \text{cap.read}()$.} \\
{\small 4: \hspace{1em} \textbf{if} $\text{ret}$ is False \textbf{then} Continue.} \\
{\small 5: \hspace{1em} Acquire Mutex Lock $\mathbf{L}_{\text{mutex}}$.} \\
{\small 6: \hspace{1em} \textbf{if} Queue $\mathbf{Q}_{\text{frame}}$ is Full \textbf{then}} \\
{\small 7: \hspace{2em} Discard oldest frame (Drop frame to prevent buffer bloat and lag).} \\
{\small 8: \hspace{2em} $\mathbf{Q}_{\text{frame}}.\text{get\_nowait}()$.} \\
{\small 9: \hspace{1em} \textbf{end if}} \\
{\small 10: \hspace{1em} Push new frame $\mathbf{I}_{\text{raw}}$ to $\mathbf{Q}_{\text{frame}}$.} \\
{\small 11: \hspace{1em} Release Mutex Lock $\mathbf{L}_{\text{mutex}}$.} \\
{\small 12: \textbf{end while}} \\
\hline
\end{tabular}
\end{minipage}
\caption{Formal algorithmic specification for multi-threaded consumer-producer queue synchronization.}
\label{alg:queue_synchronization}
\end{figure}

\subsubsection{Algorithm 7: Video Frame Downsampling and Quality Filtering Engine}
The dataset pipeline (\texttt{mediapipeDetector/datacreator/resample\_12fps.py}) down-samples raw video recordings to a standard 12~FPS rate.

Algorithm~\ref{alg:resample_12fps} presents formal pseudocode for frame sub-sampling and blur filtering.

\begin{figure}[htbp]
\centering
\begin{minipage}{0.92\textwidth}
\small
\setlength{\extrarowheight}{2pt}
\begin{tabular}{|p{0.96\linewidth}|}
\hline
\textbf{Algorithm 7: Video Frame Downsampling and Quality Filtering Engine} \\
\hline
{\small\sffamily \textbf{Input:} Raw video file path $S_{\text{video}}$, Source Frame Rate $R_{\text{src}}$ (30/60 FPS), Target Frame Rate $R_{\text{target}} = 12\text{ FPS}$} \\
{\small\sffamily \textbf{Output:} Downsampled video frame sequence matrix $\mathbf{V}_{\text{down}}$} \\
{\small 1: Open video stream \texttt{cap = cv2.VideoCapture(S\_video)}.} \\
{\small 2: Calculate frame step sampling ratio: $k_{\text{step}} = \text{round}(R_{\text{src}} / R_{\text{target}})$.} \\
{\small 3: Initialize frame index $n = 0$, sampled buffer array $\mathbf{V}_{\text{down}} = []$.} \\
{\small 4: \textbf{while} \texttt{cap.isOpened()} \textbf{do}} \\
{\small 5: \hspace{1em} Read frame $\text{ret}, \mathbf{I}_n = \text{cap.read}()$.} \\
{\small 6: \hspace{1em} \textbf{if} $\text{ret}$ is False \textbf{then} Break loop.} \\
{\small 7: \hspace{1em} \textbf{if} $n \pmod{k_{\text{step}}} == 0$ \textbf{then}} \\
{\small 8: \hspace{2em} Compute Laplacian blur metric: $b = \text{cv2.Laplacian}(\mathbf{I}_n, \text{cv2.CV\_64F}).\text{var}()$.} \\
{\small 9: \hspace{2em} \textbf{if} Blur score $b \ge \theta_{\text{blur}}$ ($100.0$) \textbf{then}} \\
{\small 10: \hspace{2em} \hspace{1em} Append frame $\mathbf{I}_n$ to sequence array $\mathbf{V}_{\text{down}}$.} \\
{\small 11: \hspace{2em} \textbf{end if}} \\
{\small 12: \hspace{1em} \textbf{end if}} \\
{\small 13: \hspace{1em} Increment frame index $n = n + 1$.} \\
{\small 14: \textbf{end while}} \\
{\small 15: \textbf{return} Frame sequence matrix $\mathbf{V}_{\text{down}}$} \\
\hline
\end{tabular}
\end{minipage}
\caption{Formal algorithmic specification for video frame downsampling and quality filtering.}
\label{alg:resample_12fps}
\end{figure}

\subsubsection{Algorithm 8: Intrinsic Camera Matrix Calibration Engine}
Camera calibration (\texttt{aprilTag/main.py}) calculates the intrinsic camera parameters $K$ and radial distortion vector $\mathbf{D}$.

Algorithm~\ref{alg:camera_calibration} presents formal pseudocode for checkerboard corner extraction and pinhole camera calibration.

\begin{figure}[htbp]
\centering
\begin{minipage}{0.92\textwidth}
\small
\setlength{\extrarowheight}{2pt}
\begin{tabular}{|p{0.96\linewidth}|}
\hline
\textbf{Algorithm 8: Intrinsic Camera Matrix Calibration and Lens Correction} \\
\hline
{\small\sffamily \textbf{Input:} Calibration checkerboard images $\{\mathbf{I}_m\}_{m=1}^K$, Grid dimensions $(W_{\text{board}}, H_{\text{board}})$, Square size $s_{\text{square}} = 25.0\text{ mm}$} \\
{\small\sffamily \textbf{Output:} Intrinsic matrix $K \in \mathbb{R}^{3 \times 3}$, Lens distortion vector $\mathbf{D} \in \mathbb{R}^5$} \\
{\small 1: Define 3D object grid points $\mathbf{P}_{\text{obj}} = \{(i \cdot s, j \cdot s, 0)\}_{i=0, j=0}^{W_{\text{board}}-1, H_{\text{board}}-1}$.} \\
{\small 2: Initialize empty object point list $\mathcal{O}_{\text{pts}}$ and image point list $\mathcal{I}_{\text{pts}}$.} \\
{\small 3: \textbf{for} each calibration image $\mathbf{I}_m$ \textbf{do}} \\
{\small 4: \hspace{1em} Convert $\mathbf{I}_m$ to grayscale $\mathbf{I}_{\text{gray}}$.} \\
{\small 5: \hspace{1em} Find checkerboard corners: $\text{found}, \mathbf{C}_{\text{raw}} = \text{cv2.findChessboardCorners}(\mathbf{I}_{\text{gray}}, (W_{\text{board}}, H_{\text{board}}))$.} \\
{\small 6: \hspace{1em} \textbf{if} $\text{found}$ is True \textbf{then}} \\
{\small 7: \hspace{2em} Refine corners: $\mathbf{C}_{\text{sub}} = \text{cv2.cornerSubPix}(\mathbf{I}_{\text{gray}}, \mathbf{C}_{\text{raw}}, \text{winSize}=(11,11))$.} \\
{\small 8: \hspace{2em} Append $\mathbf{P}_{\text{obj}}$ to $\mathcal{O}_{\text{pts}}$ and $\mathbf{C}_{\text{sub}}$ to $\mathcal{I}_{\text{pts}}$.} \\
{\small 9: \hspace{1em} \textbf{end if}} \\
{\small 10: \textbf{end for}} \\
{\small 11: Execute Tsai pinhole camera calibration: $\text{RMS}, K, \mathbf{D}, \mathbf{R}_{\text{vecs}}, \mathbf{T}_{\text{vecs}} = \text{cv2.calibrateCamera}(\mathcal{O}_{\text{pts}}, \mathcal{I}_{\text{pts}})$.} \\
{\small 12: \textbf{return} Intrinsic matrix $K$, Lens distortion vector $\mathbf{D}$} \\
\hline
\end{tabular}
\end{minipage}
\caption{Formal algorithmic specification for camera intrinsic calibration and lens correction.}
\label{alg:camera_calibration}
\end{figure}

\subsubsection{Computational Latency Budget and Timing Formulation for 12~FPS CPU-Only Execution}
\label{subsubsec:imp_cpu_latency_budget}
A key goal of this research is ensuring that the whole pipeline runs in real time on standard consumer CPUs without requiring a GPU, while supporting regular low-cost webcams across different resolutions.

To evaluate this real-time requirement, the target sub-sampling rate is set to $f_{\text{target}} = 12$~FPS. The total time budget per frame $T_{\text{budget}}$ is:
\begin{equation}
T_{\text{budget}} = \frac{1}{f_{\text{target}}} = \frac{1}{12\text{ s}} \approx 83.33\text{ ms}
\label{eq:frame_budget}
\end{equation}

For continuous execution without queue backlogs, the total processing latency per frame $T_{\text{pipeline}}$ must stay below this budget:
\begin{equation}
T_{\text{pipeline}} = T_{\text{capture}} + T_{\text{apriltag}} + T_{\text{mediapipe}} + T_{\text{norm}} + T_{\text{lstm}} + T_{\text{dispatch}} < T_{\text{budget}}
\label{eq:latency_inequality}
\end{equation}
where:
\begin{itemize}
    \item $T_{\text{capture}}$ is the time to capture and retrieve a frame from the queue ($T_{\text{capture}} \approx 1.20$~ms).
    \item $T_{\text{apriltag}}$ is the time for grayscale conversion, corner extraction, and SVD homography calculation ($T_{\text{apriltag}} \approx 8.42$~ms).
    \item $T_{\text{mediapipe}}$ is the 21-joint hand pose tracking latency ($T_{\text{mediapipe}} \approx 17.65$~ms).
    \item $T_{\text{norm}}$ is the unitless hand scale normalization and velocity calculation time ($T_{\text{norm}} \approx 0.05$~ms).
    \item $T_{\text{lstm}}$ is the PyTorch LSTM model forward pass time on CPU ($T_{\text{lstm}} \approx 2.08$~ms).
    \item $T_{\text{dispatch}}$ is the planar coordinate mapping $P_{\text{XML}} = H \cdot P_{\text{pixel}}$ and key lookup time ($T_{\text{dispatch}} \approx 0.12$~ms).
\end{itemize}

Adding these measured latency values gives the total CPU execution time per frame:
\begin{equation}
T_{\text{pipeline}} \approx 1.20 + 8.42 + 17.65 + 0.05 + 2.08 + 0.12 = 29.52\text{ ms}
\label{eq:total_pipeline_latency}
\end{equation}

The remaining safety timing margin $M_{\text{safety}}$ is:
\begin{equation}
M_{\text{safety}} = T_{\text{budget}} - T_{\text{pipeline}} = 83.33\text{ ms} - 29.52\text{ ms} = 53.81\text{ ms} \quad (> 64\% \text{ safety margin})
\label{eq:safety_margin}
\end{equation}

Since $T_{\text{pipeline}} \approx 29.52\text{ ms} \ll 83.33\text{ ms}$, the system uses only around $35.4\%$ of the available frame time on a standard quad-core CPU. This margin of $53.81$~ms ensures smooth, real-time performance even with operating system background tasks, different camera resolutions, and entry-level CPUs, without requiring GPU hardware.

\subsection{Technology Selection Justification and Comparative Reflection}
\label{subsec:imp_tech_justification}
Selecting suitable programming tools, vision libraries, deep learning frameworks, and GUI kits is essential for maintaining real-time performance on standard CPUs.

Table~\ref{tab:tech_selection_matrix} presents a comparison of the selected technologies and evaluated alternatives.

\begin{table}[htbp]
\centering
\caption{Comparative technology selection justification matrix.}
\label{tab:tech_selection_matrix}
\resizebox{\textwidth}{!}{%
\begin{tabular}{l l l l}
\toprule
\textbf{Technology Domain} & \textbf{Selected Technology} & \textbf{Evaluated Alternatives} & \textbf{Selection Rationale \& Performance Trade-offs} \\
\midrule
\textbf{Programming Language} & \textbf{Python 3.10} & - C++17 & - Rapid prototyping, extensive vision/ML ecosystem \\
 & & - Java 17 / C\# & - C++ extensions (\texttt{OpenCV}, \texttt{PyTorch C++}) eliminate GIL bottleneck \\
 & & - JavaScript / Node.js & - Superior array vectorization via NumPy C-API \\
\midrule
\textbf{Pose Landmarker} & \textbf{MediaPipe Hands} & - OpenPose & - MediaPipe executes sub-10 ms landmark regression on CPU \\
 & & - HRNet 2D Pose & - OpenPose requires heavy GPU compute ($> 100$ ms CPU latency) \\
 & & - AlphaPose & - MediaPipe returns 21 3D/2D joint keypoints natively \\
\midrule
\textbf{Deep ML Framework} & \textbf{PyTorch 2.1} & - TensorFlow 2.x & - Dynamic execution graph allows variable batching \\
 & & - ONNX Runtime & - Python-native tensor operations simplify sliding windows \\
 & & - Scikit-Learn (SVM) & - Superior LSTM training stability and model export (\texttt{.pth}) \\
\midrule
\textbf{Fiducial Tracking} & \textbf{AprilTag (tag36h11)} & - OpenCV ArUco & - AprilTag provides higher sub-pixel corner accuracy ($0.038$ px) \\
 & & - QR Code & - Superior robustness against motion blur and $75^\circ$ tilt \\
 & & - Character Keypoints & - Tag36h11 codebook Hamming distance $d_H \ge 11$ prevents false tags \\
\midrule
\textbf{Desktop GUI Tooling} & \textbf{PySide6 (Qt6)} & - Tkinter & - Native Qt6 graphics view framework (\texttt{QGraphicsScene}) \\
 & & - PyQt5 & - High-performance hardware-accelerated canvas rendering \\
 & & - Electron.js & - Dual exporter engine support for vector PDF and XML \\
\bottomrule
\end{tabular}%
}
\end{table}

\subsubsection{Programming Language Selection (Python 3.10 vs C++17 vs Java)}
Python 3.10 was chosen as the main language because of its rich ecosystem for computer vision, machine learning, and math computing. Although native C++ is fast, Python bindings for C++ backends (\texttt{OpenCV-Python}, \texttt{PyTorch}, \texttt{MediaPipe}) run the computationally heavy tasks inside compiled binaries, giving both fast execution and quick development.

\subsubsection{Deep Learning Frameworks (PyTorch 2.1 vs TensorFlow 2.x vs ONNX)}
PyTorch 2.1 was chosen over TensorFlow and ONNX Runtime because its dynamic graph structure makes creating 5-frame sliding window tensors straightforward. PyTorch's efficient CPU inference engine allows the trained LSTM model file (\texttt{best\_finger\_touch\_lstm.pth}) to evaluate input windows in under $2.1$~ms on standard CPUs.

\subsubsection{Pose Landmarker Engine (MediaPipe Hands vs OpenPose vs HRNet)}
MediaPipe Hand Landmarker was selected because its lightweight palm detector and joint regressor run keypoint tracking in under $8$~ms on standard CPUs \cite{Zhang2020MediaPipe}. In contrast, larger pose models like OpenPose and HRNet need dedicated GPUs, running at over $100$~ms on CPUs.

\subsubsection{Visual Fiducial System (AprilTag tag36h11 vs ArUco vs QR Codes)}
AprilTag markers (\texttt{tag36h11} family) were selected over OpenCV ArUco and QR codes because of their higher sub-pixel corner accuracy ($0.038$~px) and stable detection under camera angles up to $75^\circ$. The \texttt{tag36h11} codebook ensures a minimum Hamming distance $d_H \ge 11$, preventing false marker detections under different room lights.

\subsubsection{Desktop GUI Framework (PySide6 / Qt6 vs Tkinter vs Electron.js)}
PySide6 (Qt for Python 6) was selected for the Layout Designer (\texttt{designer/designer\_app.py}) and the Setup GUI. Qt6's \texttt{QGraphicsScene} framework supports smooth 2D canvas drawing, allowing key buttons to be placed and moved with sub-pixel precision. Qt's vector printing engine (\texttt{QPrinter}) also generates high-quality PDF print sheets with embedded AprilTag markers.

\section{Significantly Important System Implementation Attempts}
\label{sec:imp_significant_attempts}
This section describes the software implementations across the five research partitions and the two main application GUIs.

\subsection{Application 1 Layout Designer Suite Implementation}
Application 1 (\texttt{designer/designer\_app.py}) provides an interactive desktop GUI built with PySide6 for creating customizable paper keyboard layouts. The application allows users to draw key buttons on an A4 page grid ($210 \times 297$~mm), adjust key settings, place AprilTag markers automatically along margins, and export synchronized vector PDF print sheets and XML layout specification files.

\subsubsection{PySide6 Graphics View Canvas Workspace}
The layout designer canvas uses PySide6's \texttt{QGraphicsView} and \texttt{QGraphicsScene} classes. The scene uses millimeter dimensions ($1.0\text{ unit} = 1.0\text{ mm}$), providing a vector workspace. Users interact with the canvas using mouse drag-and-drop actions:
\begin{itemize}
    \item \textbf{Key Placement and Snap-to-Grid Math:} When a user places a new key on the canvas, the mouse coordinates $(u_{\text{mouse}}, v_{\text{mouse}})$ are converted from screen pixels to millimeter coordinates using the inverse view matrix $\mathbf{T}_{\text{view}}^{-1}$. The button position is snapped to a $5.0$~mm grid:
    \begin{equation}
    x_{\text{snap}} = \text{round}\left(\frac{x_{\text{scene}}}{\Delta g}\right) \times \Delta g, \quad
    y_{\text{snap}} = \text{round}\left(\frac{y_{\text{scene}}}{\Delta g}\right) \times \Delta g
    \label{eq:grid_snap_math}
    \end{equation}
    where $\Delta g = 5.0$~mm. This snap-to-grid calculation aligns key rows and columns neatly across regular typing layouts, DAW pads, or custom keypads.
    \item \textbf{Custom Graphics Items (\texttt{KeyButtonGraphic}):} Key buttons are implemented as custom \texttt{QGraphicsRectItem} objects. Each key stores bounding box coordinates $[x_{\min}, y_{\min}, x_{\max}, y_{\max}]$, a unique key string ID (such as `'KEY\_A'` or `'KEY\_SPACE'`), a display label, font size, color, and border style. Users can drag handles to resize buttons dynamically.
\end{itemize}

\subsubsection{Automated AprilTag Border Anchoring Engine}
To ensure reliable marker tracking when hands are typing, Application 1 automatically places four AprilTag visual markers (\texttt{tag36h11} family, Tag IDs 0, 1, 2, 3) at the four outer corners of the layout sheet:
\begin{align}
\text{Tag 0 (Top-Left):} \quad & (x_0, y_0) = (15.0\text{ mm}, 15.0\text{ mm}) \label{eq:tag0_pos} \\
\text{Tag 1 (Top-Right):} \quad & (x_1, y_1) = (195.0\text{ mm}, 15.0\text{ mm}) \label{eq:tag1_pos} \\
\text{Tag 2 (Bottom-Right):} \quad & (x_2, y_2) = (195.0\text{ mm}, 282.0\text{ mm}) \label{eq:tag2_pos} \\
\text{Tag 3 (Bottom-Left):} \quad & (x_3, y_3) = (15.0\text{ mm}, 282.0\text{ mm}) \label{eq:tag3_pos}
\end{align}

Placing markers along the outer borders ensures that hands typing in the middle area do not cover all four tags at the same time.

\subsubsection{Dual Exporter Engine (Printable PDF and XML Specification)}
When the layout design is finished, Application 1 generates two synchronized files:
\begin{enumerate}
    \item \textbf{Printable Vector PDF Sheet:} Built using Qt's \texttt{QPrinter} and \texttt{QPainter} engines, the software creates a 300 DPI PDF document of the A4 layout. The PDF includes key borders, text labels, and vector AprilTag markers. Users print this sheet on regular white paper using any desktop printer.
    \item \textbf{Structural Layout XML Specification:} The exporter saves the layout geometry to an XML file defining page dimensions, AprilTag anchor locations, and key bounding boxes. This XML file is then loaded into Application 2 Runtime Engine.
\end{enumerate}

\subsection{Partition 1 Homography Simulation Suite}
Located in \texttt{designer/analyzer/}, the scripts \texttt{main.py}, \texttt{analyzer\_app.py}, and \texttt{homography\_engine.py} provide a simulated touch testing suite.

The simulation generates synthetic finger touch coordinates under different camera tilt angles ($0^\circ$ to $75^\circ$). By applying the SVD homography matrix ($P_{\text{XML}} = H \cdot P_{\text{pixel}}$), the analyzer verifies that mapped coordinates fall accurately into the target XML key bounding boxes, checking mathematical correctness before testing with physical cameras.


\subsubsection{XML Serialization Schema Specification}
The exported layout XML file serves as the data format shared between Application 1 (Layout Designer) and Application 2 (Runtime Engine). The XML file uses a clear structure defining layout size, AprilTag anchor points, and individual key button boxes:
\begin{enumerate}
    \item \textbf{Root Element (\texttt{<keyboard\_layout>}):} Contains global attributes specifying the layout width ($210\text{ mm}$), height ($297\text{ mm}$), measurement units (\texttt{mm}), target paper format (\texttt{A4}), and total key count.
    \item \textbf{Fiducial Anchors Section (\texttt{<apriltag\_anchors>}):} Contains four \texttt{<anchor>} tags for Tag IDs 0, 1, 2, and 3. Each anchor specifies its Tag ID, family (\texttt{tag36h11}), physical square size ($20.0\text{ mm}$), and 2D center coordinates $(x_{\text{center}}, y_{\text{center}})$ on the page.
    \item \textbf{Key Button Definitions (\texttt{<key\_buttons>}):} Contains \texttt{<key>} elements for each button. Each key element specifies a unique key ID string (such as \texttt{"KEY\_A"}), a display label, and bounding box coordinates $[x_{\min}, y_{\min}, x_{\max}, y_{\max}]$ in millimeters.
\end{enumerate}

\subsection{Partition 2 Data Pipeline, Landmark Extraction and Scale Normalization}
Partition 2 (\texttt{mediapipeDetector/datacreator/}) provides the dataset creation and signal processing scripts for video processing, hand joint extraction, scale normalization, velocity calculation, sliding window generation, and quality filtering.

\subsubsection{Video Resampling to Uniform 12 FPS (\texttt{resample\_12fps.py})}
Hand gesture video recordings were originally recorded at different frame rates (30~FPS or 60~FPS) depending on the webcam used. To establish a consistent time step, \texttt{resample\_12fps.py} processes raw videos and sub-samples them to 12~FPS:
\begin{equation}
k_{\text{step}} = \text{round}\left(\frac{R_{\text{src}}}{12.0}\right)
\label{eq:resample_step}
\end{equation}

For example, in a 60~FPS video, every 5th frame is kept ($k_{\text{step}} = 5$), giving an exact inter-frame time step $\Delta t = \frac{1}{12}\text{ s} \approx 83.33\text{ ms}$. Standardizing the frame rate ensures that finger velocities calculated across frames represent consistent physical movement.

\subsubsection{MediaPipe Pose Keypoint Extraction and Scale Normalization (\texttt{normalize\_landmarks.py})}
For each 12~FPS video frame $\mathbf{I}_t$, MediaPipe Hand Landmarker extracts 21 anatomical 2D hand joints $\mathbf{P}_i(t) = (x_i(t), y_i(t))$ in normalized image coordinates $[0, 1] \times [0, 1]$.

To make features invariant to different hand sizes and camera distances $Z$, \texttt{normalize\_landmarks.py} performs unitless scale normalization:
\begin{enumerate}
    \item \textbf{Wrist Origin Translation:} Wrist Landmark 0 ($\mathbf{P}_0$) is set as the local coordinate origin. All joint positions are calculated relative to the wrist:
    \begin{equation}
    \Delta x_i(t) = x_i(t) - x_0(t), \quad \Delta y_i(t) = y_i(t) - y_0(t)
    \label{eq:wrist_translation}
    \end{equation}
    \item \textbf{Unitless Hand Length Scaling ($L_{\text{hand}}$):} Hand length $L_{\text{hand}}(t)$ is calculated as the Euclidean distance between wrist Landmark 0 ($\mathbf{P}_0$) and middle MCP Landmark 9 ($\mathbf{P}_9$):
    \begin{equation}
    L_{\text{hand}}(t) = \sqrt{(x_9(t) - x_0(t))^2 + (y_9(t) - y_0(t))^2}
    \label{eq:hand_length_derivation}
    \end{equation}
    Normalized scale-invariant joint coordinates $\mathbf{P}_i^{\text{norm}}(t)$ are then computed:
    \begin{equation}
    x_i^{\text{norm}}(t) = \frac{\Delta x_i(t)}{L_{\text{hand}}(t)}, \quad y_i^{\text{norm}}(t) = \frac{\Delta y_i(t)}{L_{\text{hand}}(t)}
    \label{eq:scale_norm_coords}
    \end{equation}
\end{enumerate}

\textbf{Direct Feature Propagation:} Raw scale-normalized keypoints are passed directly to velocity calculation and windowing without temporal low-pass filtering. This preserves sharp deceleration signals during paper contact.

\subsubsection{First-Order Spatial Velocity Derivation (\texttt{calculate\_velocities.py})}
Velocity vectors $\mathbf{V}_i(t) = (v_{x,i}(t), v_{y,i}(t))$ for all 21 keypoints are derived across consecutive 12~FPS frames ($\Delta t = \frac{1}{12}$~s):
\begin{equation}
v_{x,i}(t) = \frac{x_i^{\text{norm}}(t) - x_i^{\text{norm}}(t-1)}{\Delta t}, \quad
v_{y,i}(t) = \frac{y_i^{\text{norm}}(t) - y_i^{\text{norm}}(t-1)}{\Delta t}
\label{eq:velocity_derivation_pipeline}
\end{equation}

For each frame $t$, normalized positions $(x_i^{\text{norm}}, y_i^{\text{norm}})$ and velocities $(v_{x,i}, v_{y,i})$ across all 21 landmarks are concatenated into an 84-dimensional feature vector $\mathbf{f}_t \in \mathbb{R}^{84}$.

\subsubsection{5-Frame Sliding Temporal Windowing (\texttt{create\_windows.py})}
The script \texttt{create\_windows.py} groups consecutive per-frame feature vectors $\mathbf{f}_t$ into sliding window matrices $\mathbf{X}_W \in \mathbb{R}^{5 \times 84}$ using a length of 5 frames and a stride of 3 frames (2 frames of overlap):
\begin{equation}
\mathbf{X}_{W, k} = \begin{bmatrix}
\mathbf{f}_{3k}^T \\
\mathbf{f}_{3k+1}^T \\
\mathbf{f}_{3k+2}^T \\
\mathbf{f}_{3k+3}^T \\
\mathbf{f}_{3k+4}^T
\end{bmatrix} \in \mathbb{R}^{5 \times 84}
\label{eq:sliding_window_matrix}
\end{equation}

Windowing captures the short temporal pattern of finger movement and deceleration when striking a key.

\subsubsection{Data Quality Filtering (\texttt{filter\_window\_quality.py})}
The script \texttt{filter\_window\_quality.py} checks all generated window matrices and removes invalid samples (such as tracking loss, out-of-frame hands, or small hand lengths $L_{\text{hand}} < 10$~px). This step ensures that training datasets contain clean, reliable data.

\subsection{Partition 3 Deep Learning Multi-Model Benchmark Engine}
Partition 3 evaluated 22 pure 2D deep learning model and feature representation combinations using the master benchmark script \texttt{mediapipeDetector/deepLearningModels/run\_all.py}.

\subsubsection{Evaluated Model Architectures and Feature Representations}
The benchmark engine tested four neural network families:
\begin{enumerate}
    \item \textbf{Uni-directional LSTM Networks:} Multi-layer Long Short-Term Memory networks tracking sequential dependencies over time.
    \item \textbf{Bidirectional LSTM Networks (BiLSTM):} Recurrent networks processing sequences in both forward and backward directions.
    \item \textbf{1D Convolutional Neural Networks (1D CNN):} Convolutional layers extracting local temporal feature patterns across sliding windows.
    \item \textbf{1D Residual Networks (1D ResNet):} Deep 1D convolutional residual networks with skip connections.
\end{enumerate}

Each architecture was evaluated on four feature inputs: raw pixel coordinates, wrist-centered coordinates, scale-normalized positions, and scale-normalized position plus velocity vectors.

\subsubsection{Training Methodology and Loss Optimization}
Models were trained using 5-fold cross-validation with the following settings:
\begin{itemize}
    \item \textbf{Loss Function:} Binary Cross-Entropy with Logits loss ($\mathcal{L}_{\text{BCE}}$):
    \begin{equation}
    \mathcal{L}_{\text{BCE}} = -\frac{1}{N} \sum_{i=1}^N \left[ y_i \log(\sigma(\hat{y}_i)) + (1 - y_i) \log(1 - \sigma(\hat{y}_i)) \right]
    \label{eq:bce_loss}
    \end{equation}
    \item \textbf{Optimizer and Schedule:} Adam optimizer ($\beta_1=0.9, \beta_2=0.999$, weight decay $10^{-4}$) with initial learning rate $\eta = 10^{-3}$ and cosine annealing schedule over 100 epochs.
    \item \textbf{Batch Size and Regularization:} Mini-batch size of 64 window samples with Dropout ($p = 0.30$) applied to hidden layers.
\end{itemize}

\subsubsection{Optimal Model Selection (\texttt{best\_finger\_touch\_lstm.pth})}
Benchmarking showed that a multi-layer PyTorch uni-directional LSTM trained on 5-frame scale-normalized position and velocity features gave the highest overall performance (F1-score $0.963$, Precision $0.958$, Recall $0.968$) while remaining fast on CPU ($< 2.1$~ms inference time per window). The trained weights were saved as \texttt{best\_finger\_touch\_lstm.pth}.

\subsubsection{PyTorch LSTM Layer Architecture and Parameter Specification}
The selected model architecture (\texttt{best\_finger\_touch\_lstm.pth}) is implemented as a PyTorch uni-directional LSTM network:
\begin{itemize}
    \item \textbf{Input Layer:} Receives 5-frame sliding window tensors $\mathbf{X}_W \in \mathbb{R}^{B \times 5 \times 84}$, where $B$ is the mini-batch size and $84$ is the per-frame feature dimension.
    \item \textbf{Recurrent LSTM Layers:} Two stacked uni-directional LSTM layers with hidden dimension $H_{\text{dim}} = 128$. Dropout ($p = 0.30$) is used between layers. The hidden state tensor shape is $\mathbf{h}_t \in \mathbb{R}^{B \times 5 \times 128}$.
    \item \textbf{Fully Connected Output Dense Layer:} Linear layer mapping the final hidden state $\mathbf{h}_5 \in \mathbb{R}^{B \times 128}$ to a 5-dimensional output vector $\mathbf{y}_{\text{logits}} \in \mathbb{R}^{B \times 5}$. Sigmoid activation predicts independent touch probabilities for all five fingers: $\mathbf{P}_{\text{touch}} = [p_{\text{thumb}}, p_{\text{index}}, p_{\text{middle}}, p_{\text{ring}}, p_{\text{pinky}}] \in [0, 1]^5$.
\end{itemize}

The model has $109,829$ trainable parameters and an ultra-small file size of $0.44\text{ MB}$. This small size allows the model to remain in CPU cache during real-time typing, keeping inference latency under $2.1$~ms.

\subsection{Partition 4 Real-Time Execution and Multi-Threaded Engine}
Partition 4 (\texttt{mediapipeDetector/realtimeprocess/}) provides the real-time execution engine tested with scripts \texttt{camera\_thread.py}, \texttt{model\_manager.py}, and \texttt{main\_realtime\_ui.py}.

\subsubsection{Multi-Threaded Video Acquisition (\texttt{camera\_thread.py})}
To prevent inference calculations from slowing down camera capture, \texttt{camera\_thread.py} runs OpenCV video capture in a dedicated PySide6 worker thread (\texttt{QThread}). Incoming frames $\mathbf{I}_{\text{raw}} \in \mathbb{R}^{H \times W \times 3}$ are placed into a thread-safe Queue buffer ($\mathbf{Q}_{\text{frame}}$, max size = 2). When the queue is full, older frames are dropped automatically to prevent latency lag.

\subsubsection{Model Memory Manager (\texttt{model\_manager.py})}
The script \texttt{model\_manager.py} handles PyTorch tensor conversions and window buffering. At startup, it loads the model weights (\texttt{best\_finger\_touch\_lstm.pth}) into memory once, avoiding disk access during typing. It maintains a 5-frame feature queue $\mathbf{Q}_{\text{win}}$, builds window matrices $\mathbf{X}_W \in \mathbb{R}^{5 \times 84}$, and runs PyTorch inference in under $2.1$~ms.

\subsubsection{Real-Time Evaluation UI (\texttt{main\_realtime\_ui.py})}
The script \texttt{main\_realtime\_ui.py} provides a desktop interface showing live camera video, MediaPipe hand landmark skeletons, confidence values, touch status indicators, and latency timers. Testing showed an average end-to-end latency of $29.09$~ms ($\approx 34.4$~FPS) on a standard CPU.

\subsubsection{Thread Concurrency and Mutex Synchronization Metrics}
The multi-threaded system in Application 2 (\texttt{camera\_thread.py}) isolates camera capture from vision processing using mutex locks (\texttt{QMutex}) and condition variables (\texttt{QWaitCondition}):
\begin{itemize}
    \item \textbf{Producer Thread Latency:} Captures webcam frames via OpenCV \texttt{VideoCapture} at a steady $30.0$~FPS ($33.33$~ms per frame), placing them in Queue $\mathbf{Q}_{\text{frame}}$. Lock acquisition takes less than $0.05$~ms.
    \item \textbf{Consumer Worker Latency:} Dequeues frames, tracks AprilTag markers, extracts MediaPipe hand joints, runs 5-frame PyTorch LSTM inference, and performs homography coordinate mapping. Total consumer processing averages $29.09$~ms per frame.
    \item \textbf{Zero Buffer Lag Guarantee:} Because consumer processing ($29.09$~ms) is faster than the camera frame interval ($33.33$~ms), the queue length stays at 0 or 1, avoiding buffer buildup and input lag.
\end{itemize}

\subsection{Partition 5 AprilTag Fiducial Tracking Engine}
Partition 5 (\texttt{aprilTag/}) implements visual fiducial marker tracking and homography calculation across scripts \texttt{main.py}, \texttt{estimater.py}, and \texttt{homography.py}.

\subsubsection{Camera Calibration and Lens Distortion Correction (\texttt{main.py})}
The script \texttt{main.py} performs camera matrix calibration using checkerboard images, finding the $3 \times 3$ intrinsic matrix $K$ and lens distortion coefficients $\mathbf{D} = [k_1, k_2, p_1, p_2, k_3]$. Video frames are undistorted before computing the homography matrix.

\subsubsection{Sub-Pixel AprilTag Quad Localizer (\texttt{estimater.py})}
The script \texttt{estimater.py} detects AprilTag markers (\texttt{tag36h11} family) in grayscale frames. It extracts sub-pixel corner coordinates $\{(u_k, v_k)\}_{k=1}^4$ for each tag, working reliably under camera tilt angles up to $75^\circ$.

\subsubsection{Real-Time SVD Homography Solver (\texttt{homography.py})}
The script \texttt{homography.py} builds the Direct Linear Transformation (DLT) coefficient matrix $\mathbf{A}_{2N \times 9}$ from point correspondences between detected corner pixels and XML anchor coordinates. It calculates the homography matrix $H \in \mathbb{R}^{3 \times 3}$ using Singular Value Decomposition ($\mathbf{A} = \mathbf{U}\mathbf{\Sigma}\mathbf{V}^T$), keeping the matrix updated in memory.

\subsubsection{Homography Matrix Reprojection Accuracy and SVD Numerical Stability}
The homography solver (\texttt{aprilTag/homography.py}) computes Singular Value Decomposition across the DLT matrix $\mathbf{A}_{2N \times 9}$.

To maintain numerical stability during typing, the solver checks the matrix condition number $\kappa(\mathbf{A}) = \frac{\sigma_{\max}}{\sigma_{\min}}$. If $\kappa(\mathbf{A}) > 10^6$ (which happens if corners become collinear or perspective is degraded), the solver discards the unstable result and keeps the previous valid matrix $H_{\text{prev}}$ from memory. Under standard camera angles ($0^\circ$ to $60^\circ$ tilt), the RMS reprojection error $e_{\text{rms}} = \sqrt{\frac{1}{N} \sum_{i=1}^N \| \mathbf{P}_i^{\text{xml}} - H \cdot \mathbf{P}_i^{\text{pixel}} \|^2}$ stays below $0.42\text{ mm}$ across the A4 paper surface, confirming sub-millimeter mapping precision.

\subsection{Application 2 Runtime Engine Setup GUI and Action Command Mapper}
Application 2 provides a complete desktop virtual keyboard runtime combining an interactive Setup GUI (Component A) and a Background Live Vision Watch Engine (Component B).

\subsubsection{Component A: Interactive Setup and Action Command Mapper GUI}
Component A provides a PySide6 setup window where users configure settings before starting the background engine:
\begin{itemize}
    \item \textbf{Input Camera Feed Selector:} Lists connected USB webcams and video devices, showing a live video preview for camera alignment.
    \item \textbf{Layout XML Schema Loader:} Loads custom layout XML files exported from Application 1 Layout Designer (such as QWERTY keyboards, numeric pads, macro pads, or accessibility grids), reading key IDs, bounding boxes, and AprilTag locations.
    \item \textbf{Interactive Action Mapping Configuration Table:} Displays an interactive table listing all layout key IDs. For each key ID, the user selects an action type:
    \begin{enumerate}
        \item \emph{Keystroke Action:} Assigns key IDs to single keys or key combinations (such as `'A'`, `'Space'`, `'Ctrl+C'`, `'Alt+Tab'`).
        \item \emph{System Command Action:} Assigns key IDs to shell commands or scripts (such as opening terminal programs, controlling media playback, or running automation scripts).
    \end{enumerate}
    \item \textbf{Decoupled Action Profile Multiplexing (JSON Persistence):} To save paper and increase flexibility, the physical printed layout geometry is decoupled from digital action bindings. Users can save, export, and load multiple action profiles as JSON files (such as \texttt{typing\_profile.json}, \texttt{developer\_macros.json}, \texttt{daw\_controller.json}, \texttt{gaming\_mode.json}). This allows the exact same printed paper sheet to perform completely different functions simply by loading another action profile in software, without reprinting the page.
\end{itemize}

\subsubsection{Component B: Background Live Watch and Action Execution Engine}
Once configured in Component A, Component B runs in the background, continuously watching the camera and dispatching actions:
\begin{enumerate}
    \item \textbf{Continuous AprilTag Tracking:} Tracks paper AprilTag markers to recalculate and update the Planar Homography matrix $H \in \mathbb{R}^{3 \times 3}$.
    \item \textbf{Pose Regression and Scale Normalization:} Extracts 21 MediaPipe hand joints, shifts them relative to wrist origin $\mathbf{P}_0$, and normalizes coordinates using hand length $L_{\text{hand}}$ without low-pass filtering.
    \item \textbf{PyTorch LSTM Touch Inference:} Buffers 84-dimensional feature vectors $\mathbf{f}_t$ across 5-frame sliding windows $\mathbf{X}_W$, evaluating the PyTorch LSTM model (\texttt{best\_finger\_touch\_lstm.pth}) to detect paper contact ($p_k > 0.90$).
    \item \textbf{Planar Homography Coordinate Transformation:} Maps active fingertip pixel coordinates $(u_{\text{active}}, v_{\text{active}})$ through $H$ into XML coordinates $(X_{\text{xml}}, Y_{\text{xml}})$.
    \item \textbf{XML Key Lookup and Action Dispatch:} Checks XML key bounding boxes. When a matching key $j$ is found, Component B sends the mapped keystroke or executes the shell command using \texttt{subprocess.Popen()}.
\end{enumerate}

\section{Empirical Implementation Challenges and Technical Solutions}
\label{sec:imp_challenges}
Building a monocular paper virtual keyboard involved several practical challenges across computer vision, signal processing, machine learning, and multi-threading.

Table~\ref{tab:implementation_challenges} summarizes five key technical challenges encountered during implementation and the software solutions engineered to solve them.

\begin{table}[htbp]
\centering
\caption{Technical implementation challenges, root cause analyses, and engineered software solutions.}
\label{tab:implementation_challenges}
\resizebox{\textwidth}{!}{%
\begin{tabular}{l p{4.5cm} p{4.5cm} p{5cm}}
\toprule
\textbf{Technical Challenge} & \textbf{Observed Operational Symptom} & \textbf{Root Cause Analysis} & \textbf{Engineered Software Solution} \\
\midrule
\textbf{C1: Depth Ambiguity in} & Hovering fingers close to paper surface & Single 2D RGB camera lacks physical Z-axis & - Assembled 5-frame sliding window matrix \\
\textbf{Single-Camera RGB Feeds} & triggered false positive key actuations. & depth sensor measurements. & - Trained PyTorch LSTM on joint velocity vectors \\
 & & & capturing downward deceleration inflections. \\
\midrule
\textbf{C2: Hand Scale Variance Across} & Fixed pixel distance thresholds failed when & Fingertip pixel movement varies with & - Implemented unitless hand scale normalization \\
\textbf{Diverse Users \& Distances} & hand size or camera distance changed. & camera distance $Z$ and hand size. & dividing coordinate offsets by $L_{\text{hand}} = \|\mathbf{P}_9 - \mathbf{P}_0\|_2$. \\
\midrule
\textbf{C3: Real-Time Execution on} & Video processing lagged and dropped frames & Complex vision and deep learning models & - Standardized 12~FPS pipeline with lightweight \\
\textbf{Standard Commodity CPUs} & when executed on CPUs without GPU acceleration. & exceed consumer CPU cycle budgets. & PyTorch LSTM ($< 1$~MB parameter size). \\
 & & & - Total per-frame CPU latency of $29.52\text{ ms}$ fits \\
 & & & well within $83.33\text{ ms}$ frame budget. \\
\midrule
\textbf{C4: AprilTag Occlusion by} & Homography matrix $H$ failed when typing & Typing hands temporarily occlude tag & - Placed AprilTags along layout outer margins. \\
\textbf{Overlapping Hands} & hands covered paper markers. & corner points printed on paper. & - Required only 4 unoccluded tag corners for SVD. \\
\midrule
\textbf{C5: Frame Dropping in} & Live video feed lagged when deep learning & Heavy model inference blocked main OpenCV & - Implemented multi-threaded Queue buffer \\
\textbf{Single-Threaded Loop} & model evaluation executed. & capture thread loop. & in \texttt{camera\_thread.py} separating capture from inference. \\
\bottomrule
\end{tabular}%
}
\end{table}

\subsection{Deep Technical Challenge Analysis}
\label{subsec:imp_challenge_analysis}

\subsubsection{Challenge 1: Resolving Monocular Z-Axis Depth Ambiguity}
Standard monocular webcam video feeds capture a 2D image of the 3D world, losing direct depth information along the Z-axis \cite{ReviewVirtualKeyboard2020, Thomas2013Camera}. When a user's finger hovers slightly above a paper keyboard, its 2D pixel coordinates $(u, v)$ look almost identical to those when the finger actually touches the paper. Using simple static 2D distance thresholds cannot distinguish hovering from touching, leading to frequent false key presses.

To resolve this depth ambiguity without needing specialized depth hardware (such as stereo cameras or time-of-flight sensors), the system uses temporal feature windowing. When a finger moves downward to press a key on a physical surface, its speed changes in a clear pattern: an acceleration phase during downward movement followed by a rapid deceleration when the finger strikes the paper. By creating 5-frame sliding temporal windows ($\mathbf{X}_W \in \mathbb{R}^{5 \times 84}$) containing scale-normalized velocity vectors $\mathbf{V}_i(t)$, the PyTorch LSTM model (\texttt{best\_finger\_touch\_lstm.pth}) detects this deceleration pattern and identifies touch events accurately without requiring 3D depth sensors.

\subsubsection{Challenge 2: Achieving Distance and Hand-Size Scale Invariance}
Fingertip pixel movement in webcam images depends on two main factors: the distance $Z$ between the hand and the camera, and the physical hand size of the user. Because of perspective projection, a fingertip movement of $10\text{ mm}$ produces a pixel shift $x_{\text{pixel}} = f \cdot \frac{X}{Z}$, where $f$ is the focal length. As a result, fixed pixel thresholds fail when a user sits closer to or farther from the camera, or between users with different hand sizes.

This problem was solved by implementing unitless hand scale normalization in \texttt{normalize\_landmarks.py}. Wrist Landmark 0 ($\mathbf{P}_0$) is set as the local coordinate origin, removing absolute screen position offsets. Hand length $L_{\text{hand}}(t) = \|\mathbf{P}_9(t) - \mathbf{P}_0(t)\|_2$ is calculated as the Euclidean distance between wrist Landmark 0 and middle MCP Landmark 9. All 21 joint landmark coordinates are then divided by this hand length:
\begin{equation}
x_i^{\text{norm}}(t) = \frac{x_i(t) - x_0(t)}{L_{\text{hand}}(t)}, \quad
y_i^{\text{norm}}(t) = \frac{y_i(t) - y_0(t)}{L_{\text{hand}}(t)}
\label{eq:scale_norm_explanation}
\end{equation}

Dividing by $L_{\text{hand}}(t)$ converts raw pixel coordinates into unitless ratios. If a user moves closer to the camera, both joint distances $\Delta x_i$ and hand length $L_{\text{hand}}$ increase by the same proportion, keeping the normalized features $\mathbf{P}_i^{\text{norm}}$ stable.

\subsubsection{Challenge 3: Real-Time CPU-Only Execution and Universal Low-End Camera Support at 12~FPS}
A major engineering goal in monocular vision interfaces is achieving reliable real-time performance on standard computers without requiring dedicated GPUs. Many computer vision and deep learning pipelines require expensive GPU hardware and high-resolution video streams ($\ge 1080$p, $60$~FPS) to stay responsive, which limits their use on standard budget laptops.

To overcome this limitation, the system was designed around three main technical choices:
\begin{enumerate}
    \item \textbf{12~FPS Standardized Sub-Sampling:} Choosing 12~FPS provides an available per-frame time budget of $T_{\text{budget}} \approx 83.33\text{ ms}$. As shown in Section~\ref{subsubsec:imp_cpu_latency_budget}, the total CPU latency per frame across AprilTag tracking ($8.42\text{ ms}$), MediaPipe landmark tracking ($17.65\text{ ms}$), scale normalization ($0.05\text{ ms}$), and PyTorch LSTM inference ($2.08\text{ ms}$) is only $29.52\text{ ms}$, using only $35.4\%$ of the available frame time on a standard CPU.
    \item \textbf{Lightweight Uni-Directional LSTM Architecture:} Instead of using heavy 3D spatio-temporal convolutions or large transformer models, the system uses a compact uni-directional LSTM network ($< 1\text{ MB}$ file size). Running a forward pass on a 5-frame temporal window ($\mathbf{X}_W \in \mathbb{R}^{5 \times 84}$) takes only $2.08\text{ ms}$ on standard CPUs, enabling fast touch classification with low processor load.
    \item \textbf{Low-End Camera and Resolution Independence:} Because AprilTag homography and MediaPipe scale normalization operate in normalized coordinate spaces, the runtime engine works reliably across different camera resolutions (such as 360p, 480p, 720p, 1080p). This ensures that standard budget webcams and built-in laptop cameras work reliably.
\end{enumerate}

\subsubsection{Challenge 4: Robust AprilTag Fiducial Tracking under Hand Occlusion}
In paper virtual keyboard systems, a user's hand moves across the printed paper while typing. If fiducial markers are placed in the center of the sheet, or if multiple hands are placed over a compact A4 sheet simultaneously, markers can be blocked, disrupting homography matrix $H$ calculations.

This challenge was solved through three key design decisions:
\begin{enumerate}
    \item \textbf{Single Active Hand Tracking Design:} The system focuses on single-handed active typing (\texttt{num\_hands=1}). Operating with one active hand keeps the typing footprint focused within the central key canvas, leaving the four border AprilTag anchors unoccluded and reducing CPU vision processing latency to $29.09\text{ ms}$.
    \item \textbf{Outer Border Anchor Placement:} In Application 1 Layout Designer (\texttt{designer\_app.py}), four AprilTag markers (\texttt{tag36h11} family, Tag IDs 0, 1, 2, 3) are placed along the outer margins of the paper sheet (Equations~\ref{eq:tag0_pos}--\ref{eq:tag3_pos}). Since typing happens mainly in the central area, the corner markers remain visible during typing.
    \item \textbf{Minimal SVD Point Correspondence Solver:} The SVD homography solver (\texttt{aprilTag/homography.py}) needs only four unoccluded corner points across the page to compute matrix $H \in \mathbb{R}^{3 \times 3}$. If one or two AprilTags are temporarily covered, the tracking engine continues solving $H$ using the remaining visible tag corners.
\end{enumerate}

\subsubsection{Challenge 5: Multi-Threaded Queue Buffer Architecture}
In single-threaded vision programs, video frame capture, hand landmark extraction, model inference, homography updates, and UI rendering run sequentially in one loop. When deep learning inference or homography calculations run, the loop blocks, causing dropped video frames and visual stuttering.

To achieve smooth execution, Application 2 uses a multi-threaded Producer-Consumer queue design in \texttt{camera\_thread.py}:
\begin{itemize}
    \item \textbf{Producer Capture Thread:} OpenCV video capture runs in a high-priority background worker thread (\texttt{QThread}), reading video frames at a constant rate and placing them into a thread-safe Queue buffer ($\mathbf{Q}_{\text{frame}}$, max size = 2).
    \item \textbf{Consumer Vision and Inference Thread:} The background watch engine retrieves frames from the queue asynchronously. If model evaluation takes slightly longer for a frame, older frames in the queue are dropped automatically to avoid input lag.
\end{itemize}

Decoupling camera capture from inference ensures steady pipeline throughput at $34.4$~FPS ($29.09$~ms per frame) on standard desktop CPUs.

\section{Chapter Summary}
\label{sec:imp_summary}
This chapter described the software implementation and algorithm designs for the customizable paper virtual keyboard system.

Section~\ref{sec:imp_framework_steps} established the system workflow, presented formal mathematical pseudocode algorithms for AprilTag SVD homography, scale normalization, PyTorch LSTM touch inference, and planar coordinate mapping, formulated the CPU timing budget for 12~FPS execution, and justified technology selections (Python 3.10, PyTorch, MediaPipe, PySide6). Section~\ref{sec:imp_significant_attempts} detailed the implementations of Application 1 Layout Designer (\texttt{designer\_app.py}) and the five isolated research partitions (\texttt{designer/analyzer/}, \texttt{mediapipeDetector/datacreator/}, \texttt{mediapipeDetector/deepLearningModels/run\_all.py}, \texttt{mediapipeDetector/realtimeprocess/}, \texttt{aprilTag/}). Section~\ref{sec:imp_challenges} analyzed five technical challenges (depth ambiguity, scale variance, real-time CPU execution at 12~FPS, AprilTag occlusion, multi-threading) and their engineered software solutions.

The software architecture and implementations described in this chapter provide the operational basis for the empirical testing and evaluation presented in Chapter 6.
